\documentclass[10pt,a4paper]{amsart}
\usepackage[margin=19mm]{geometry}
\usepackage{amsmath,amssymb,mathtools}
\usepackage[scr=rsfs,cal=euler]{mathalfa}
\usepackage{xfrac}
\usepackage[unicode,hidelinks,bookmarks=false]{hyperref}
\newcommand{\ie}{i.e.\ }
\newcommand{\cf}{cf.\ }
\newcommand{\resp}{respectively\ }
\newcommand{\Subsection}{\subsection}
\providecommand{\nobreakdash}{\nobreak\hbox{-}}
\setlength{\emergencystretch}{2em}
\multlinegap0pt
\usepackage{xcolor}
\newtheorem{theorem}{Theorem}
\def\bbR{\mathbb{R}}
\title{long version here}
\author{Olivia Dumitrescu, Rick Miranda}
\date{Source: arXiv:2104.14141v4 (CC BY 4.0)}
\begin{document}
\maketitle

\noindent\emph{Source and licence.} long version here. Version arXiv:2104.14141v4 (CC BY 4.0), distributed by arXiv under the Creative Commons Attribution 4.0 International licence, http://creativecommons.org/licenses/by/4.0/, as recorded in the arXiv licence metadata for this exact version and shown on the abstract page https://arxiv.org/abs/2104.14141v4. Full licence notice: \texttt{licenses/arxiv-2104.14141-CC-BY-4.0.txt}.\par\smallskip

\noindent\textit{Editorial scope.} Counted unit: the article's theorem cone (source0.tex lines 3254--3259). The article's complete original proof of it is delivered with it (source0.tex lines 3261--3297, one \texttt{proof} environment of 37 lines). No mathematics is written, repaired or re-derived: the statement and the proof are the article's own lines, in the article's own notation, and no retained line is substituted or re-flowed.  Retained uncounted as the source-local context the proof needs: lines 3228--3230.  Nothing was omitted from any retained span, so the package carries no omission record.  No retained line contains a citation, so the wrapper carries no bibliography and no bibliography entry.  No retained line contains a figure environment, an image-inclusion command or drawing code, so no image is delivered and the package declares no asset dependency.  Editorial formatting only, in addition: an AMS \texttt{amsart} preamble that restates the article's own declarations and macros used by the retained lines, namely the environments \texttt{theorem} on separate counters, together with the standard packages those lines need.  The retained environments print in this document as Theorem 1 in order of appearance, whereas the article prints the same items with the numbering of its own full document.  The complete native source of the article as distributed in the arXiv e-print for this exact version is bundled unchanged as \texttt{sources/110859/source0.tex}.
\begin{equation}\label{equation1}
\mathcal{E}:=\{\text{(i)-Weyl line}  \;\text{ for }\; i \in \{0,1\}\} \subset A^{r-1}(Y).
\end{equation}

\begin{theorem}\label{effective cone}
 If $s \leq r+3$ then
 $$\text{Eff}_{\mathbb{R}} Y^r_s = \mathcal{E}_{\geq 0}.$$
In other words, $(0)$- and $(1)$-Weyl lines give the extremal rays
for the cone of movable curves in $\mathbb{P}^r$ with $r+3$ points blown up.
\end{theorem}

\begin{proof}
We will first compute the set $\mathcal{E} \subset A^{r-1}(Y)$ of Equation \eqref{equation1}.

It is easy to see that $(0)$-Weyl lines consists only in a pencil of lines through one point and the rational normal curves of degree $r$ passing through $r+2$ general points.
\[
\alpha_{i}=h-e_i \text{  and  } \alpha'_{i}=nh - \sum_{j\neq i} e_j.
\]
The $(0)$-Weyl lines $\alpha_{i}$ give the facets $(A_{i})$.

We will leave to the reader to check that the $(1)$-curves are of the form
\[
\beta_{t, I(t)} = (t+1)r-t - \sum_{i\in I(t)} (t+1)e_i -  \sum_{i\notin I(t)} t e_i
\]
for every $-1\leq t\leq l+\alpha$ and $|I(t)|=r-2t+1$.
Indeed, ordering the multiplicities in a decreasing order,
one can see that performing a standard Cremona transformation to the last $r+1$ points, then
\[
Cr(\beta_{t, I(t)}) = \beta_{t+1, I(t+1)}.
\]

The $(1)$-Weyl line of minimal degree correspond to $t=-1$,
i.e. the general hyperplane class $\beta_{-1, I(-1)}=h$,
while $(1)$-Weyl line of maximal degree correspond to $t$ a half of $r$
i.e. when
for $\alpha=0$ ($r$ is even) is the quasihomogeneous curve $\beta_{l, I(l)}$,
and if $\alpha=1$, a homogeneous curve $\beta_{l+1, I(l+1)}$.
The $(1)$-Weyl line $\beta_{t, I(t)}$ with $|I(t)|=r-2t+1$ give facets $(B_{r,I(t)})$.
We obtain
\begin{equation}
\begin{split}
 \mathcal{E}&=\{ \{\alpha_i, \alpha'_i\}_{1\leq i\leq r+3}, \text{ and } \{\beta_{t, I(t)}\}_{ -1 \leq t \leq l} \},\\
\mathcal{E}_{\geq 0}&=\{D\in \text{Pic}(Y)_{\bbR} \text{ so that } D \cdot C\geq 0\ \text{for } C\in \mathcal{E}\},\\
&=\{D\in \text{Pic}(Y)_{\bbR} \text{ satisfying } (A_{i})\ \text{ and } (B_{n,I(t)})\},\\
&= \text{Eff}_{\mathbb{R}} Y.
\end{split}
\end{equation}
\end{proof}

\end{document}
